\section{Introduction}
\label{sec:intro}

Type inference with subtyping collects the typing requirements of a program as
\emph{subtype constraints} $s\le t$ between type expressions with unknowns
\cite{Mitchell91,FuhMishra90,AikenWimmers93,EifrigSmithTrifonov95,Pottier98}.
Types are finite or infinite trees. When the type language has a least type
$\Bot$ and a greatest type $\Top$, the subtype order is \emph{non-structural}:
$\Bot\le t\le\Top$ holds for every tree $t$, so the order relates trees of
different shapes. For instance, with a binary constructor $f$ that is covariant
in both arguments,
\[
  f(\Bot,\,f(\Bot,\Bot))\;\le\;f(f(\Top,\Bot),\,\Top),
\]
although the two trees have no common shape.

Three questions about a \emph{constraint} $\Gamma$, a finite set of such
inequalities read as their conjunction, have been studied for this order. \emph{Satisfiability} asks whether $\Gamma$ has a solution; it
is decidable in polynomial time
\cite{KozenPalsbergSchwartzbach94,PalsbergWandOKeefe97}. \emph{Entailment}
asks whether every solution of $\Gamma$ satisfies a further inequality
$x\le y$; it is the basis of constraint simplification
\cite{Pottier96,TrifonovSmith96,Rehof97,Pottier01}. Its decidability, posed
by Rehof in 1996 as Problem~16 of the TLCA list of open problems
\cite{Rehof96TLCA}, was studied by Niehren and Priesnitz with automata
\cite{NiehrenPriesnitz99,NiehrenPriesnitz03} and settled positively only
recently \cite{Zhao26}. Finally, the full \emph{first-order theory} of the
order is undecidable \cite{SuAikenNiehrenPriesnitzTreinen02}, and for a
binary covariant constructor, with or without $\Top$, this already holds for
sentences with the quantifier prefix $\exists\forall\exists$
\cite{Zhao26U}.

Satisfiability is an existential ($\Sigma_1$) sentence
$\exists\bar x.\,\bigwedge\Gamma$ in which no negation occurs, and entailment
is a universal ($\Pi_1$) sentence
$\forall\bar x.\,\neg\bigwedge\Gamma\vee x\le y$ of a special form: its
quantifier-free part has a single positive inequality. This paper decides the
whole first level of the quantifier hierarchy.

\begin{theorem}[Main theorem]\label{thm:main}
Let $f$ be a constructor of arity $n\ge0$, each of whose arguments is covariant
or contravariant. It is decidable whether a Boolean combination of $\Sigma_1$
and $\Pi_1$ sentences over subtype constraints with the signature $\BFT$ is
true, over arbitrary trees, over regular trees, and over finite trees. A
sentence is true over the regular trees exactly when it is true over all
trees. The same holds for the signature $\BF$ when $f$ is covariant. The
decision procedures are uniform in the arity and the variance.
\end{theorem}

The function type constructor is the case $n=2$ with a contravariant first
argument. Since a $\Pi_1$ sentence is the negation of a $\Sigma_1$ sentence,
the sentences of the theorem are the Boolean combinations of $\Sigma_1$
sentences; we call the set of the true ones the \emph{$\Sigma_1$ theory}.
Truth over finite trees can differ from truth over all trees: for a unary
constructor, $\exists x.\,x\le f(x)\wedge f(x)\le x$ is true over regular
trees, with the infinite tree $f(f(\cdots))$ as witness, and false over finite
trees.

\subparagraph*{Disjunctive entailment.}
A $\Sigma_1$ sentence is true exactly when some clause of a disjunctive normal
form of its quantifier-free part has a solution. After every compound subterm
is named by a variable, such a clause has the form
\begin{equation}\label{eq:intro-clause}
  \Gamma\;\wedge\;\neg(x_1\le y_1)\;\wedge\;\dots\;\wedge\;\neg(x_k\le y_k),
\end{equation}
where $\Gamma$ is a constraint, a conjunction of inequalities, and
$x_1,y_1,\dots,x_k,y_k$ are variables. It has a solution exactly when
$\Gamma$ does \emph{not} entail the disjunction
$x_1\le y_1\vee\dots\vee x_k\le y_k$. So the theorem is equivalent to the
decidability of \emph{disjunctive entailment}
\[
  \Gamma\;\entails\;x_1\le y_1\vee\dots\vee x_k\le y_k,
\]
uniformly in the number $k$ of \emph{queries}. Plain entailment is the case
$k=1$ and satisfiability the case $k=0$.

Disjunctive entailment does not reduce to plain entailment in any direct way,
because a disjunction can be entailed although none of its disjuncts is. For
the function type constructor every tree is $\Bot$ or lies above the least
function type $f(\Top,\Bot)$, so $\forall x.\;x\le\Bot\vee f(\Top,\Bot)\le
x$ is true, but neither disjunct holds for all $x$. A countermodel of a
disjunction must refute all disjuncts with a single solution
(\cref{sec:sent-nonreduction}).

\subparagraph*{From countermodels to holes.}
A failed inequality $x\le y$ has a \emph{witness}: a path at which the labels
of the values of $x$ and $y$ are in the wrong order. For $k$ queries we obtain
$k$ witness paths $w_1,\dots,w_k$, which are words over the argument positions
$\{1,\dots,n\}$. We show that a tuple of words is the tuple of witnesses of
some countermodel exactly when it is a \emph{hole} of a finite automaton
compiled from $\Gamma$ and the queries, a \emph{structural reader}
(\cref{sec:readers}). The reader has one start state for each query. It reads
$w_i$ from the $i$th start state, and the tuple is a hole if every run stays
alive and ends in a \emph{rejected} state, and no \emph{admitted} pair of
states occurs at a \emph{comparison}: a pair of cuts, one in a word $w_i$ and
one in a word $w_j$, such that the rest of $w_j$ after its cut is a prefix of
the rest of $w_i$ after its cut. The words $w_i$ and $w_j$ may be the same
word or different ones. The pairs of states at the comparisons between $w_i$
and $w_j$ form the \emph{channel} between the two words. A comparison is a
place where a subtree that one witness forces into the solution meets a
subtree that another witness must keep out of it. Here a state is rejected
when the reader has rejected every way in which the constraint could refute
the witness, so rejected states are the good ones.

Whether a reader has a hole is not a regular property of the tuple of words,
since comparisons relate unboundedly long suffixes at unboundedly distant
cuts. Decidability follows from a bound on the length of a shortest hole. We
write $\expo_0(x)=x$ and $\expo_{h+1}(x)=2^{\expo_h(x)}$ for iterated
exponentials.

\begin{theorem}[Small holes]\label{thm:intro-small}
There is a computable function $\Hole$ such that every structural reader with
$N$ states and $k$ start states that has a hole has a hole of total length at
most $\Hole(k,N)$, whatever the size of its alphabet, and
\[
  \Hole(k,N)\;\le\;\expo_4\bigl((k+N+20)^{2k+6}\bigr).
\]
\end{theorem}

For a fixed number $k$ of queries the bound is four iterated exponentials
of a polynomial in $N$. A reader compiled from a constraint with $v$
variables has at most $2^{4v}+1$ states, so the witnesses of a countermodel can be chosen of length
at most five iterated exponentials of a polynomial in the size of the
sentence, and the $\Sigma_1$ theory is decidable in elementary space
(\cref{sec:complexity}). The problem is PSPACE-hard, because plain entailment
is \cite{HengleinRehof98}. The procedures establish decidability and an
elementary upper bound; they are not practical algorithms.

\subparagraph*{Why several queries are harder than one.}
The decision procedure for plain entailment \cite{Zhao26} also reduces a
countermodel to a hole and bounds a shortest hole, for readers with one start
state. Its argument does not extend to several words, for three reasons, which
are summarized in \cref{tab:intro-compare}.

\emph{Several witnesses in one solution.} With $k$ queries a hole is a tuple
of words, and there are comparisons between different words. Each
query can fail in two ways, at a path where the smaller side is too large or
at a path where the larger side is too small, so there is one reader for each
of the $2^k$ vectors of \emph{modes}.

\emph{Rejected tails.} A run enters the rejected states at some cut and then
stays in them. With one word, no admitted pair can involve a cut of the
\emph{rejected tail}, the part of the word after that cut, so a hole can be
cut off shortly after it. With several words, a cut in the rejected tail of
one word can still be the host of an admitted comparison with another word.
If the run can be continued forever, the tail may have to be long in order to
avoid the other words, and its length must be bounded independently of their
lengths. This is the largest part of the proof (\cref{sec:tails}). We replace
such tails by ultimately periodic words, \emph{rays}; we show that rays can be
chosen so that the constraints between their lengths form a solvable system
of difference constraints; and we compress the periodic stretches of the
resulting hole without creating an admitted pair on any channel.

\emph{Shortening all words together.} Deleting a block from one word of a
tuple can create a comparison with another word. The shortening of
\cref{sec:fresh} therefore follows one hierarchy of \emph{markers} for the
whole tuple and deletes the same repeated blocks from all words at the same
time.

\begin{table}[t]
\caption{Plain entailment and disjunctive entailment.}
\label{tab:intro-compare}
\small\setlength{\tabcolsep}{4pt}
\begin{tabular}{@{}>{\raggedright\arraybackslash}p{0.21\linewidth}>{\raggedright\arraybackslash}p{0.33\linewidth}>{\raggedright\arraybackslash}p{0.40\linewidth}@{}}
\toprule
 & Plain entailment \cite{Zhao26} & Disjunctive entailment (this paper) \\
\midrule
witness of a countermodel & one path & one path per query, all realized by one solution \\
\addlinespace
readers & one per failure side, with one start state & one per mode vector, $2^k$ in all, with $k$ start states \\
\addlinespace
comparisons & within one word & within and between the words of a tuple \\
\addlinespace
after the first rejection & no admitted pair; short tails & tails can still be hosts; bounded by rays and compression \\
\addlinespace
shortening & hierarchy of one word & one marker hierarchy, the same blocks deleted in all words \\
\bottomrule
\end{tabular}
\end{table}

\subparagraph*{Contributions.}
\Cref{fig:overview} shows the structure of the proof.
\begin{itemize}
\item We prove that the $\Sigma_1$ theory of non-structural subtyping is
  decidable for one constructor of any arity and variance with $\Bot$ and
  $\Top$, over arbitrary, regular and finite trees, and for a covariant
  constructor with $\Bot$ only (\cref{thm:main}, proved in
  \cref{sec:decision}).
\item We give an exact reduction of the countermodels of a disjunctive
  entailment to the holes of structural readers with $k$ start states
  (\cref{thm:reduction}).
\item We prove the small-hole theorem for structural readers
  (\cref{thm:small-hole}) from two results of independent interest in
  combinatorics on words with automata: a bound on rejected tails
  (\cref{thm:tails}) and a bound on minimal tuples of fresh words
  (\cref{thm:fresh}).
\item We show that the bound is elementary (\cref{thm:elementary}), discuss
  the gap to the PSPACE lower bound, and give families of sentences whose
  shortest witnesses are long (\cref{sec:complexity}).
\item The decision procedures, their correctness, the small-hole theorem and
  its elementary bound are mechanized in Lean~4 with Mathlib
  (\cref{sec:mechanization}).
\end{itemize}

%% overview -- the proof chain.
\begin{figure}[t]
\centering
\begin{tikzpicture}[
  font=\footnotesize,
  box/.style={draw, rounded corners=2pt, align=center, inner sep=3.5pt, minimum height=8mm},
  thm/.style={box, fill=black!8},
  arr/.style={-{Stealth[length=2mm]}, semithick},
  lbl/.style={font=\scriptsize, align=center}]
%% first row: the reductions
\node[box] (sent) at (0,0) {Boolean combination of\\ $\Sigma_1$ and $\Pi_1$ sentences};
\node[box] (disj) at (5.0,0) {disjunctive entailment\\ $\Gamma\entails x_1\le y_1\vee\dots\vee x_k\le y_k$};
\node[box] (hole) at (10.1,0) {hole of one of $2^k$\\ structural readers};
\draw[arr] (sent) -- node[lbl, above] {\S\ref{sec:sentences}} (disj);
\draw[arr] (disj) -- node[lbl, above] {\S\ref{sec:readers}} (hole);
%% second row: the small-hole theorem and its two halves
\node[thm] (small) at (10.1,-2.0) {small-hole theorem\\ (\cref{thm:small-hole})};
\draw[arr, densely dashed] (small) -- node[lbl, left] {decides by bounded\\ search (\S\ref{sec:decision-arbitrary})} (hole);
\node[thm] (cover) at (6.5,-2.0) {cover of a hole\\ (\cref{lem:cover})};
\draw[arr] (cover) -- (small);
\node[thm] (tail) at (1.6,-1.45) {rejected tails (\S\ref{sec:tails}): rays,\\ difference constraints, compression};
\node[thm] (fresh) at (1.6,-2.75) {fresh tuples (\S\ref{sec:fresh}):\\ marker hierarchy, block deletion};
\draw[arr] (tail.east) -- ++(0.5,0) |- ([yshift=1.5mm]cover.west);
\draw[arr] (fresh.east) -- ++(0.75,0) |- ([yshift=-1.5mm]cover.west);
%% third row: classes of trees
\node[box] (cls) at (5.0,1.55) {disjunctive entailment over regular trees, over finite trees,\\
  and for the signature $\BF$};
\draw[arr] (cls) -- node[lbl, right] {\S\ref{sec:decision-regular}--\ref{sec:decision-bf}} (disj);
\end{tikzpicture}
\caption{The proof chain. Solid arrows between white boxes are reductions:
a sentence is reduced to disjunctive entailments (naming of subterms), and a
failed entailment over arbitrary trees with the signature $\BFT$ to a hole
of a finite reader (witnesses and modes); the other classes of trees and the
signature $\BF$ are reduced to this case. The gray boxes are the small-hole
theorem and the results it follows from. It bounds the length of a shortest
hole, so a bounded search decides the entailment (dashed arrow). Its proof
bounds the rejected tails of a hole and then the hole itself, read as a
tuple of fresh words in a larger automaton, its cover.}
\label{fig:overview}
\end{figure}


\Cref{sec:prelim} fixes trees, the order and sentences. \Cref{sec:sentences}
reduces sentences to disjunctive entailment, and \cref{sec:readers} reduces
disjunctive entailment to holes. \Cref{sec:overview} states the two bounds
and composes them. \Cref{sec:fresh,sec:tails} explain their proofs, with the
longer arguments in the appendices. \Cref{sec:decision} gives the decision
procedures, including regular and finite trees and the signature $\BF$, and
\cref{sec:complexity} the complexity bounds. \Cref{sec:mechanization}
describes the mechanization, \cref{sec:related} discusses related work, and
\cref{sec:future,sec:conclusion} discuss future work and conclude. The
notation is collected in \cref{app:notation}.
